\documentclass[final]{beamer}

\usepackage[size=custom,width=120,height=90,scale=1.0]{beamerposter}
\usepackage{lmodern}
\usepackage{booktabs}
\usepackage{array}
\usepackage{graphicx}
\usepackage{ragged2e}
\usepackage{multicol}
\usepackage{enumitem}
\usepackage{tikz}
\usepackage{ifthen}

\setbeamertemplate{navigation symbols}{}
\setlength{\parskip}{0.35em}
\setlength{\parindent}{0pt}

\definecolor{MainBlue}{RGB}{24,62,114}
\definecolor{AccentTeal}{RGB}{32,125,122}
\definecolor{SoftGray}{RGB}{245,247,250}
\definecolor{WarmRed}{RGB}{155,55,55}

\setbeamercolor{background canvas}{bg=white}
\setbeamercolor{block title}{fg=white,bg=MainBlue}
\setbeamercolor{block body}{fg=black,bg=SoftGray}

\newcommand{\smartimage}[2]{%
  \IfFileExists{#1}{%
    \includegraphics[width=#2]{#1}%
  }{%
    \fbox{%
      \begin{minipage}[c][0.18\textheight][c]{#2}
        \centering
        \small\ttfamily Missing image\\
        \scriptsize #1
      \end{minipage}%
    }%
  }%
}

\newcommand{\caseimage}[2]{%
  \IfFileExists{#1}{%
    \includegraphics[width=#2]{#1}%
  }{%
    \fbox{%
      \begin{minipage}[c][0.16\textheight][c]{#2}
        \centering
        \small\ttfamily #1
      \end{minipage}%
    }%
  }%
}

\title{Escaping Local Minima in End-to-End AMR Navigation with Stall-Triggered Policy Switching}
\author{Jiabao Cao, Shengqin Jiang, Wenjing Chen}
\institute{Faculty of Science and Engineering, University of Nottingham Ningbo China}

\begin{document}
\begin{frame}[t]
\begin{columns}[t,totalwidth=\textwidth]

% ---------------- Column 1 ----------------
\begin{column}{0.31\textwidth}
  \begin{block}{Problem and Baseline}
    \justifying
    We study end-to-end navigation for an AMR in a simulator. The baseline is CNNTD3, which performs well in standard scenes but fails in structured hard scenes such as U-traps.

    \vspace{0.6em}
    \textbf{Main question:} how can the robot recover from local minima without losing precision in narrow passages?

    \vspace{0.6em}
    \textbf{Baseline result:}
    \begin{itemize}[leftmargin=1.2em]
      \item Standard environment SR: 87\%
      \item U-trap SR: 0\%
      \item Narrow door SR: 100\%
    \end{itemize}

    \vspace{0.4em}
    % Suggested image: baseline failure trajectory
    % Replace with a real U-trap failure screenshot if available.
    \smartimage{docs/img/smoke_test.png}{0.95\linewidth}
  \end{block}

  \begin{block}{Failure Taxonomy}
    \justifying
    \begin{tabular}{>{\raggedright\arraybackslash}p{0.3\linewidth} p{0.63\linewidth}}
      \toprule
      Type & Example / meaning \\
      \midrule
      Local minimum & U-trap: goal points through wall \\
      Precision loss & Exploration policy fails narrow door \\
      Oscillation & Repeated small turns or back-and-forth motion \\
      Config mismatch & NeuPAN in compact 10 x 10 m scenes \\
      \bottomrule
    \end{tabular}
  \end{block}

  \begin{block}{What Did Not Work}
    \begin{itemize}[leftmargin=1.2em]
      \item Naive reward shaping made the policy too cautious.
      \item A single exploration-heavy policy improved traps but hurt narrow doors.
      \item STPS v3 added complexity without clear gain over v2.
      \item NeuPAN was not convincing in the compact benchmark setting.
    \end{itemize}
  \end{block}
\end{column}

% ---------------- Column 2 ----------------
\begin{column}{0.37\textwidth}
  \begin{block}{Method: STPS}
    \justifying
    \textbf{STPS = Stall-Triggered Policy Switching.}

    Instead of forcing one policy to do everything, the robot uses a precision policy by default. When it detects stall or oscillation, it switches to an exploration policy for a short escape phase, then switches back after progress.

    \vspace{0.8em}
    \centering
    \begin{tikzpicture}[node distance=1.4cm, every node/.style={draw, rounded corners, align=center, minimum width=0.78\linewidth, minimum height=1cm, font=\small}]
      \node (main) {Precision policy};
      \node (detect) [below of=main] {Detect stall / oscillation};
      \node (escape) [below of=detect] {Switch to exploration policy\\escape for 120 steps};
      \node (recover) [below of=escape] {Switch back after progress};
      \node (main2) [below of=recover] {Precision policy};
      \draw[->, thick] (main) -- (detect);
      \draw[->, thick] (detect) -- (escape);
      \draw[->, thick] (escape) -- (recover);
      \draw[->, thick] (recover) -- (main2);
    \end{tikzpicture}

    \vspace{0.8em}
    \textbf{Trigger settings:}
    \begin{itemize}[leftmargin=1.2em]
      \item Stall: 20-step displacement $< 0.15$ m
      \item Oscillation: at least 5 direction reversals in 12 steps
      \item Recovery: displacement $> 0.5$ m
    \end{itemize}
  \end{block}

  \begin{block}{Ablation Summary}
    \centering
    \begin{tabular}{lcc}
      \toprule
      Method & U-trap & Narrow door \\
      \midrule
      CNNTD3 baseline & 0\% & 100\% \\
      Exploration policy & 100\% & 0\% \\
      STPS v2 & 75$\pm$7\% & 100\% \\
      \bottomrule
    \end{tabular}

    \vspace{0.6em}
    \justifying
    The ablation shows the core tension: exploration helps escape traps, but it can damage precise alignment. STPS resolves this by switching policies instead of merging them.
  \end{block}

  \begin{block}{Quantitative Results}
    \centering
    \small
    \begin{tabular}{lccccc}
      \toprule
      Method & Std. & U & D-U & Door & Avg. \\
      \midrule
      CNNTD3 & 87\% & 0\% & 69\% & 100\% & 67\% \\
      NeuPAN & 0\% & 0\% & 0\% & 0\% & 0\% \\
      STPS v2 & 88\% & 75\% & 100\% & 100\% & 94\% \\
      \bottomrule
    \end{tabular}
  \end{block}
\end{column}

% ---------------- Column 3 ----------------
\begin{column}{0.31\textwidth}
  \begin{block}{Case Gallery}
    \justifying
    \textbf{3 successes and 3 failures.} Use these six slots for screenshots or trajectory frames.

    \vspace{0.5em}
    \begin{tabular}{p{0.48\linewidth} p{0.48\linewidth}}
      \caseimage{docs/img/cases/success_u_trap.png}{\linewidth} &
      \caseimage{docs/img/cases/success_double_u.png}{\linewidth} \\
      \scriptsize STPS U-trap success &
      \scriptsize STPS Double-U success \\
      \caseimage{docs/img/cases/success_narrow_door.png}{\linewidth} &
      \caseimage{docs/img/cases/failure_baseline_u_trap.png}{\linewidth} \\
      \scriptsize STPS narrow-door success &
      \scriptsize CNNTD3 U-trap failure \\
      \caseimage{docs/img/cases/failure_explore_narrow_door.png}{\linewidth} &
      \caseimage{docs/img/cases/failure_neupan_compact.png}{\linewidth} \\
      \scriptsize Exploration-policy failure &
      \scriptsize NeuPAN compact-scene failure \\
    \end{tabular}
  \end{block}

  \begin{block}{Interpretation}
    \begin{itemize}[leftmargin=1.2em]
      \item The baseline is strong in simple scenes but weak in recovery.
      \item Exploration and precision are in tension.
      \item STPS keeps both abilities by switching at the right time.
    \end{itemize}
  \end{block}

  \begin{block}{Limitations}
    \begin{itemize}[leftmargin=1.2em]
      \item Thresholds are hand-designed.
      \item Evaluation is simulator-only.
      \item Exact shortest-path efficiency is not fully logged.
      \item Delay-compensation / NeuPAN work is exploratory, not the final claim.
    \end{itemize}
  \end{block}

  \begin{block}{Takeaway}
    \centering
    \Large\bfseries
    Failure analysis led to a simpler and better final method.
  \end{block}
\end{column}

\end{columns}

\vspace{0.5em}
\begin{center}
  \small
  Final submission items: source code, environment notes, exact commands, validation table, 3 success + 3 failure cases, focused improvement or ablation, short failure explanation, poster PDF.
\end{center}

\end{frame}
\end{document}

